\chapter{Metodologia}
\label{cap:metodologia}

% -----------------------------------------------------------------
% Objetivo do capítulo:
% Descrever como a pesquisa é conduzida, do recorte da aplicação
% até a estratégia de avaliação. O capítulo precisa permitir que
% um leitor técnico reproduza os passos principais.
%
% Distinção de escopo:
%   - Este capítulo descreve o MÉTODO.
%   - O Capítulo 4 descreve a IMPLEMENTAÇÃO (arquitetura
%     concreta da ferramenta, pipeline, indicadores).
%   - O Capítulo 5 traz os RESULTADOS.
% -----------------------------------------------------------------


\section{Caracterização da pesquisa}
\label{sec:caracterizacao_pesquisa}

% Conteúdo a desenvolver:
%   - Abordagem exploratória, orientada por problema e
%     objetivos, sem hipótese formal confirmatória.
%   - Pesquisa aplicada, conduzida sobre dados reais de uma
%     fábrica de fios esmaltados.
%   - Justificar a escolha por técnicas não supervisionadas
%     (ausência de rótulo de causa no nível do evento).


\section{Contexto da aplicação}
\label{sec:contexto_aplicacao}

% Conteúdo a desenvolver:
%   - Descrição da fábrica e do recorte adotado.
%   - Linhas de produção consideradas.
%   - Janela temporal dos dados analisados.


\section{Fontes de dados}
\label{sec:fontes_dados}

% Conteúdo a desenvolver:
%   - Dados operacionais das máquinas (sensores, periodicidade,
%     variáveis disponíveis).
%   - Histórico de rompimentos (registro de evento, anotação
%     manual quando existente).
%   - Dados de lotes, materiais e fornecedores.
%   - Dados de contexto produtivo (ordem de produção,
%     calibres, regulagens).
%   - Relatório de devoluções (ground truth supervisionado
%     para validação externa).


\section{Preparação e estruturação dos dados}
\label{sec:preparacao_dados}

% Conteúdo a desenvolver:
%   - Integração das bases.
%   - Limpeza e padronização.
%   - Tratamento de faltantes.
%   - Definição operacional de evento e de exposição.
%   - Cruzamento entre tabelas (lote, máquina, fornecedor,
%     janela de operação).


\section{Módulo de detecção de anomalias operacionais}
\label{sec:modulo_anomalia}

% Conteúdo a desenvolver:
%   - Variáveis analisadas (sensores selecionados, justificativa
%     técnica do conjunto de 14 sensores do Pilar EQ).
%   - Estratégia de modelagem (Isolation Forest e seus
%     parâmetros).
%   - Forma de detectar e quantificar desvio.
%   - Tratamento de combinações sem cobertura (rota SEM_MODELO
%     e fallback ao Pilar MP).


\section{Módulo de análise estatística histórica}
\label{sec:modulo_estatistico}

% Conteúdo a desenvolver:
%   - Construção de indicadores de risco por lote, material,
%     fornecedor e máquina.
%   - Cálculo de exposição operacional e taxa relativa de
%     rompimento.
%   - Regra de cruzamento entre máquinas para fortalecer ou
%     enfraquecer a suspeita de matéria-prima.


\section{Lógica de integração híbrida}
\label{sec:integracao_hibrida}

% Conteúdo a desenvolver:
%   - Como as evidências operacionais e históricas são
%     combinadas para produzir um diagnóstico único.
%   - Regras de decisão entre Pilar EQ e Pilar MP.
%   - Saída esperada do sistema (categorias de origem
%     provável e nível de confiança).
%   - Tratamento de empate e de evidência conflitante.


\section{Estratégia de avaliação}
\label{sec:estrategia_avaliacao}

% Conteúdo a desenvolver:
%   - Comparação entre módulo de anomalia isolado, módulo
%     estatístico isolado e abordagem híbrida.
%   - Avaliação por coerência técnica com conhecimento de
%     especialistas.
%   - Validação contra casos históricos documentados,
%     incluindo o caso IBRAME 482530.
%   - Validação externa contra o relatório de devoluções.
%   - Limites de cada estratégia de avaliação.
